
\documentclass{article}
\usepackage{colortbl}
\usepackage{makecell}
\usepackage{multirow}
\usepackage{supertabular}

\begin{document}

\newcounter{utterance}

\centering \large Interaction Transcript for game `wordle\_withclue', experiment `medium\_frequency\_words\_clue\_no\_critic', episode 12 with reflexion{-}Qwen3.5{-}9B.
\vspace{24pt}

{ \footnotesize  \setcounter{utterance}{1}
\setlength{\tabcolsep}{0pt}
\begin{supertabular}{c@{$\;$}|p{.15\linewidth}@{}p{.15\linewidth}p{.15\linewidth}p{.15\linewidth}p{.15\linewidth}p{.15\linewidth}}
    \# & \multicolumn{2}{c}{Player} && \multicolumn{2}{c}{Game Master} \\
    \hline

    \theutterance \stepcounter{utterance}  
    & & & \multicolumn{4}{p{0.6\linewidth}}{
        \cellcolor[rgb]{0.9,0.9,0.9}{
            \makecell[{{p{\linewidth}}}]{
                \texttt{\tiny{[P1$\langle$GM]}}
                \texttt{You have attempted this same task before and did not succeed. These are the plans you wrote after those attempts.} \\
\texttt{Trial \#1: New plan: The last attempt was aborted because the explanation referenced the word "startup," which is six letters long, violating the rule that the target word must be exactly five lowercase letters. Additionally, the explanation failed to explicitly mention how the provided clue ("like venture capital investments") would be used to deduce the specific five{-}letter word, as required by the instructions. Next time, I will ensure the explanation strictly adheres to the format by: 1) selecting a valid five{-}letter word that fits the clue (e.g., "funds" or "money"), 2) explicitly stating how the clue guides the choice of that specific five{-}letter word, and 3) confirming the word length and validity before submitting the guess.} \\
\\ 
\texttt{Trial \#2: New plan: The last attempt was aborted because the explanation contained the word "startup" (6 letters) and the word "assets" (6 letters), violating the strict rule that the target word must be exactly five lowercase letters. Additionally, the explanation failed to explicitly connect the clue "like venture capital investments" to the specific five{-}letter word chosen, as required by the instructions. Next time, I will ensure the explanation strictly adheres to the format by: 1) selecting a valid five{-}letter word that fits the clue (e.g., "money", "funds", "cashs", or "stock"), 2) explicitly stating how the clue guides the choice of that specific five{-}letter word, and 3) confirming the word length and validity before submitting the guess.} \\
\\ 
\texttt{Trial \#3: New plan: The previous attempts were aborted because the explanations contained invalid six{-}letter words ("startup", "assets") and failed to explicitly connect the clue to the chosen five{-}letter word. Next time, I will strictly adhere to the format by: 1) selecting only valid five{-}letter words that fit the clue "like venture capital investments" (e.g., "money", "funds", "stock", "cashs"), 2) explicitly stating in the explanation how the clue guides the choice of that specific five{-}letter word, and 3) ensuring no words longer than five letters appear in the explanation text before submitting the guess.} \\
\\ 
\\ 
\texttt{You are a language wizard who likes to guess words by using the given rules.} \\
\\ 
\texttt{Welcome to Wordle! You have six attempts to guess the target word, a valid English word of five lowercase letters (a{-}z). Please use the tags "explanation:" and "guess:" to provide a concise explanation for each guess.} \\
\\ 
\texttt{To help you make an informed guess, you will receive a clue for the word, such as} \\
\texttt{clue: snowy white.} \\
\\ 
\texttt{Here is an example guess based on the clue:} \\
\texttt{explanation: In the fairy tail Snow White, the girl is killed because she eats a poisoned apple. And the word apple has 5 letters.} \\
\texttt{guess: apple} \\
\\ 
\texttt{After each guess, your answer will be validated, and you will receive feedback indicating which letters are correct (green), which letters are correct but in the wrong position (yellow), and which letters are incorrect (red). This feedback can be useful in determining which letters to include or exclude in your next guess.} \\
\\ 
\texttt{For example, the feedback for "apple" might be:} \\
\texttt{guess\_feedback: a<yellow> p<yellow> p<green> l<yellow> e<red>} \\
\\ 
\texttt{The explanation should contain details about how both the guess\_feedback and clue work together in tandem to formulate a new guess.} \\
\\ 
\texttt{Do not generate any other text.} \\
\\ 
\texttt{Let's begin with your first guess.} \\
\\ 
\texttt{clue: like venture capital investments} \\
            }
        }
    }
     \\ \\

    \theutterance \stepcounter{utterance}  
    & \multicolumn{4}{p{0.6\linewidth}}{
        \cellcolor[rgb]{0.561,0.737,0.561}{
            \makecell[{{p{\linewidth}}}]{
                \texttt{\tiny{[P1$\rangle$GM]}}
                \texttt{explanation: The clue "like venture capital investments" suggests a word related to financial backing or capital. The word "funds" is a valid five{-}letter English word that directly relates to money provided for investment purposes. It fits the length requirement and the semantic meaning of the clue.} \\
\texttt{guess: funds} \\
            }
        }
    }
    & & \\ \\

    \theutterance \stepcounter{utterance}  
    & & & \multicolumn{2}{p{0.3\linewidth}}{
        \cellcolor[rgb]{0.9,0.9,0.9}{
            \makecell[{{p{\linewidth}}}]{
                \texttt{\tiny{[GM$|$GM]}}
                \texttt{attempts = 1} \\
\texttt{target\_word = risky} \\
\texttt{target\_word\_clue = like venture capital investments} \\
\texttt{guess = funds} \\
\texttt{guess\_feedback = f<red> u<red> n<red> d<red> s<yellow>} \\
            }
        }
    }
    & & \\ \\

    \theutterance \stepcounter{utterance}  
    & & & \multicolumn{4}{p{0.6\linewidth}}{
        \cellcolor[rgb]{0.9,0.9,0.9}{
            \makecell[{{p{\linewidth}}}]{
                \texttt{\tiny{[P1$\langle$GM]}}
                \texttt{guess\_feedback: f<red> u<red> n<red> d<red> s<yellow>} \\
\\ 
\texttt{Provide your response only in this format.} \\
\texttt{explanation: details} \\
\texttt{guess: word} \\
            }
        }
    }
     \\ \\

    \theutterance \stepcounter{utterance}  
    & \multicolumn{4}{p{0.6\linewidth}}{
        \cellcolor[rgb]{0.561,0.737,0.561}{
            \makecell[{{p{\linewidth}}}]{
                \texttt{\tiny{[P1$\rangle$GM]}}
                \texttt{explanation: The feedback shows that 's' is in the correct position (green) while 'f', 'u', 'n', 'd' are incorrect or misplaced (red/yellow). Since the clue "like venture capital investments" strongly implies money or capital, and 's' is confirmed at the end, the word likely ends in 's'. The previous letters 'f', 'u', 'n', 'd' are not part of the solution. Given the context of investments and the confirmed 's', the word might be "cashs" (though non{-}standard) or perhaps "funds" was a trap. However, considering common 5{-}letter words for investments ending in 's', "shares" is 6 letters. Let's reconsider: maybe the word is "money"? No, ends in y. "Capital" is 7. "Assets" is 6. "Stocks" is 6. "Funds" was the guess. Feedback: f{-}red, u{-}red, n{-}red, d{-}red, s{-}yellow. Wait, the feedback says s is yellow, not green. My previous reading was wrong. s is yellow, meaning it's in the word but not at the end. So the word contains 's' somewhere else. Let's try "stocks" but it's 6 letters. "Sales"? No. "Shares" is 6. "Assets" is 6. "Money" no s. "Cash" is 4. "Cashs" is 5 but bad grammar. Maybe "funds" is wrong. What 5 letter word for investments has an s not at the end? "Stock" has no s at end. "Bond" no s. "Note" no s. "Fund" no s. "Venture" is 7. "Capital" is 7. "Equity" is 6. "Profit" is 6. "Loss" is 4. "Gain" is 4. "Risk" is 4. "Asset" is 5. Does "asset" fit? Clue: like venture capital investments. Yes, assets are what investments buy. But "asset" has no 's' at the end. My guess was "funds". Feedback: f{-}red, u{-}red, n{-}red, d{-}red, s{-}yellow. So 's' is in the word, not at the end. 'f','u','n','d' are not in the word. So the word has 5 letters, contains 's', does not contain f,u,n,d. And relates to investments. "Stock" has s at end? No, stock ends in k. "Stock" has s at index 1. S{-}t{-}o{-}c{-}k. If I guess "stock", s is at 0. Feedback for funds: s is at 4 (yellow). So s is in the word, not} \\
            }
        }
    }
    & & \\ \\

    \theutterance \stepcounter{utterance}  
    & & & \multicolumn{2}{p{0.3\linewidth}}{
        \cellcolor[rgb]{0.9,0.9,0.9}{
            \makecell[{{p{\linewidth}}}]{
                \texttt{\tiny{[GM$|$GM]}}
                \texttt{Error: The response should contain the 'guess:' keyword exactly once.} \\
            }
        }
    }
    & & \\ \\

    \theutterance \stepcounter{utterance}  
    & & & \multicolumn{2}{p{0.3\linewidth}}{
        \cellcolor[rgb]{0.9,0.9,0.9}{
            \makecell[{{p{\linewidth}}}]{
                \texttt{\tiny{[GM$|$GM]}}
                \texttt{game\_result = ABORT} \\
            }
        }
    }
    & & \\ \\

    \theutterance \stepcounter{utterance}  
    & & \multicolumn{4}{p{0.6\linewidth}}{
        \cellcolor[rgb]{0.95,0.95,0.95}{
            \makecell[{{p{\linewidth}}}]{
                \texttt{\tiny{[P1$\leftrightarrow$P]}}
                \texttt{Post{-}episode Reflexion (not sent to GM):} \\
\texttt{New plan: The previous attempts were aborted because the explanations contained invalid six{-}letter words like "startup" and "assets," violating the strict rule that the target word must be exactly five lowercase letters. Additionally, the reasoning became convoluted and failed to explicitly connect the clue "like venture capital investments" to the specific five{-}letter word chosen. Next time, I will strictly adhere to the format by: 1) selecting only valid five{-}letter words that fit the clue (e.g., "money", "funds", "stock", "cash"), 2) explicitly stating in the explanation how the clue guides the choice of that specific five{-}letter word without mentioning any longer words, and 3) ensuring the explanation is concise and directly leads to a legal guess before submitting.} \\
            }
        }
    }
    & \\ \\

\end{supertabular}
}

\end{document}
